\begin{ouverture}[calcul]{1}{Multiplier et diviser des nombres relatifs}{Partie 1 · Nombres et calculs}
{Produits · Quotients · Enchaînements · Programmes de calcul}
\objectif{Multiplier deux nombres relatifs et justifier la règle des signes.}
\objectif{Diviser des relatifs et enchaîner les quatre opérations.}
\objectif{Nommer un calcul et traduire un programme de calcul.}
\prerequis{Addition, soustraction et droite graduée des nombres relatifs (Cinquième).}
\end{ouverture}
\pagestyle{otMain}\essentiel\fontsize{11}{15}\selectfont
\phantomsection\label{cours-1-1}
\essSection{1}{Construire le produit}
\begin{definition}Deux nombres sont opposés si leur somme vaut 0 : l’opposé de $a$ est $-a$. Soustraire un nombre revient à ajouter son opposé : $a-b=a+(-b)$. Sur une droite graduée, deux nombres opposés sont symétriques par rapport à l’origine.\end{definition}
\begin{propriete}[Un seul facteur négatif]Le produit d’un nombre positif et d’un nombre négatif est négatif. Par exemple $3\times(-5)=(-5)+(-5)+(-5)=-15$, et $(-5)\times3=3\times(-5)$ car l’ordre des facteurs ne change pas un produit.\end{propriete}
\begin{propriete}[Deux facteurs négatifs]Par distributivité, $(-7)\times(-4)+(-7)\times4=(-7)\times(-4+4)=0$. Donc $(-7)\times(-4)$ est l’opposé de $(-7)\times4=-28$ : il vaut 28. Règle des signes : le produit de deux nombres non nuls est positif si leurs signes sont identiques, négatif sinon. Multiplier par $-1$ donne l’opposé.\end{propriete}
\begin{exemple}Pour plusieurs facteurs non nuls, on compte les facteurs négatifs : $(-2)\times(-3)\times(-4)=-24$, car il y en a trois, nombre impair. Un produit qui contient un facteur nul est nul.\end{exemple}
\begin{attention}Deux signes « moins » n’ont pas toujours le même rôle. Dans $-4-3$, on soustrait 3 : le résultat est $-7$. Dans $(-4)\times(-3)$, on multiplie deux nombres négatifs : le résultat est 12.\end{attention}
\clearpage
\phantomsection\label{cours-1-2}
\essSection{2}{Diviser et enchaîner}
\begin{propriete}Pour $b\ne0$, le quotient $a:b$, noté aussi $\frac ab$, est le nombre qui, multiplié par $b$, donne $a$. Ainsi $(-12):4=-3$ car $(-3)\times4=-12$ : la règle des signes d’un quotient est celle du produit. On ne divise jamais par zéro.\end{propriete}
\begin{savoirfaire}[Une ligne par étape]Calculer d’abord les parenthèses, puis les produits et quotients de gauche à droite, enfin les sommes et différences de gauche à droite. Réécrire les parties inchangées pour conserver la structure.\end{savoirfaire}
\begin{exemple}$(-3)+(-4)\times7=(-3)+(-28)=-31$ et $-5\times(2-7)=-5\times(-5)=25$. Dans $6-2\times(-3+7)=6-8=-2$, la parenthèse est calculée d’abord.\end{exemple}
\begin{definition}[Nommer un calcul]Un calcul porte le nom de la dernière opération effectuée : $(-3)+(-4)\times7$ est une somme, $-5\times(2-7)$ est un produit. Le programme « choisir $-2$, ajouter 6, multiplier par $-3$ » s’écrit en une seule expression $(-2+6)\times(-3)$ : c’est un produit, qui vaut $-12$.\end{definition}
\begin{savoirfaire}[Contrôler]Prévoir le signe, estimer l’ordre de grandeur, puis vérifier par une opération inverse. Dans un problème, préciser le sens du signe et l’unité : $-8$ euros représente une baisse de 8 euros.\end{savoirfaire}

\clearpage\rubrique{Découverte}
\begin{blocLivre}{ex-1-1}\begin{exercice}[1]{Un produit positif}
Calculer $(-7)\times(-6)$.
\end{exercice}\end{blocLivre}
\begin{blocLivre}{ex-1-2}\begin{exercice}[1]{Un produit négatif}
Calculer $2{,}4\times(-5)$.
\end{exercice}\end{blocLivre}
\begin{blocLivre}{ex-1-3}\begin{exercice}[1]{Un quotient}
Calculer $(-54):(-9)$ et $18:(-4)$.
\end{exercice}\end{blocLivre}
\begin{blocLivre}{ex-1-4}\begin{exercice}[1]{Opposé ou produit}
Calculer $-8-(-3)$ et $(-8)\times(-3)$.
\end{exercice}\end{blocLivre}
\begin{blocLivre}{ex-1-5}\begin{exercice}[1]{Nommer puis calculer}
Dire si chaque calcul est une somme, un produit ou un quotient, puis le calculer : $A=(-3)+(-4)\times7$ ; $B=-5\times(2-7)$ ; $C=[12-(-6)]:(-3)$.
\end{exercice}\end{blocLivre}

\rubrique{Entraînement}
\begin{blocLivre}{ex-1-6}\begin{exercice}[2]{Compter les signes}
Sans effectuer les multiplications, donner le signe de $(-2)\times5\times(-3)\times(-7)$, puis calculer.
\end{exercice}\end{blocLivre}
\begin{blocLivre}{ex-1-7}\begin{exercice}[2]{Respecter les priorités}
Calculer $A=14-3\times(-2)$ et $B=(14-3)\times(-2)$.
\end{exercice}\end{blocLivre}
\begin{blocLivre}{ex-1-8}\begin{exercice}[2]{Une expression imbriquée}
Calculer $18:[-2\times(5-8)]$.
\end{exercice}\end{blocLivre}
\begin{blocLivre}{ex-1-9}\begin{exercice}[2]{Des multiplications à trou}
Compléter : $(-6)\times\ldots=42$ ; $\ldots\times(-2{,}5)=-10$ ; $7\times\ldots=-3$.
\end{exercice}\end{blocLivre}
\begin{blocLivre}{ex-1-10}\begin{exercice}[2]{Une erreur de calcul}
Nora écrit $24:(-6)\times2=24:(-12)=-2$. Corriger.
\end{exercice}\end{blocLivre}

\rubrique{Problèmes}
\begin{blocLivre}{ex-1-11}\begin{exercice}[3]{Évolution de température}
À midi, il fait $3\,^\circ$C. Pendant 5 heures, la température varie de $-1{,}6\,^\circ$C par heure. Quelle température obtient-on ?
\end{exercice}\end{blocLivre}
\begin{blocLivre}{ex-1-12}\begin{exercice}[3]{Un programme de calcul}
Programme : choisir un nombre ; le multiplier par $-3$ ; ajouter 7 ; multiplier le résultat par $-2$.\begin{questions}
\item Appliquer ce programme au nombre $-4$.
\item Écrire ce calcul en une seule expression.
\item Cette expression est-elle une somme ou un produit ?
\end{questions}
\end{exercice}\end{blocLivre}
\begin{blocLivre}{ex-1-13}\begin{exercice}[3]{Un nombre inconnu}
Quel nombre faut-il multiplier par $-2{,}5$ pour obtenir 15 ? Vérifier.
\end{exercice}\end{blocLivre}
\begin{blocLivre}{ex-1-14}\begin{exercice}[3]{La machine à signes}
Le produit de 19 nombres négatifs non nuls est-il positif ? Que devient son signe si l’on remplace un facteur par son opposé, puis si l’on remplace un facteur par zéro ?
\end{exercice}\end{blocLivre}
\begin{blocLivre}{ex-1-15}\begin{exercice}[3]{Un calcul à contrôler}
Calculer $C=-3[4-2(-5+8)]+12:(-4)$. Donner un contrôle du signe final.
\end{exercice}\end{blocLivre}

\clearpage
